\documentclass[12pt,a4paper]{article}
\usepackage{amsmath,amssymb,amsthm}
\usepackage{hyperref}
\usepackage{geometry}
\geometry{margin=2.5cm}
\usepackage{enumitem}
\usepackage{booktabs}

\newtheorem{theorem}{Theorem}[section]
\newtheorem{lemma}[theorem]{Lemma}
\newtheorem{corollary}[theorem]{Corollary}
\newtheorem{proposition}[theorem]{Proposition}
\theoremstyle{definition}
\newtheorem{definition}[theorem]{Definition}
\theoremstyle{remark}
\newtheorem{remark}[theorem]{Remark}

\title{The Ontology of Entanglement in Recognition Science:\\
Shared Ledger Constraints and Bell Violations}

\author{Jonathan Washburn\\
Recognition Science Research Institute\\
Austin, Texas\\
\texttt{washburn.jonathan@gmail.com}}

\date{March 2026}

\begin{document}
\maketitle

\begin{abstract}
We derive quantum entanglement from the Recognition Science (RS)
ledger structure.  Entanglement is not a mysterious ``spooky
action at a distance'' but a shared constraint on the ledger:
two subsystems that share a common ledger entry are entangled.
The inseparability theorem---that the joint state of entangled
subsystems cannot be factored---follows from the constraint that a
single ledger entry binds both outcomes.  Bell inequality violations
arise because the shared-entry correlations
$E(a, b) = -\cos(a - b)$ exceed what any local hidden-variable
model can produce.  The Tsirelson bound
$|S| \leq 2\sqrt{2} \approx 2.83$ is the maximum correlation
achievable from a shared binary-outcome entry.  Entanglement does
not transmit information faster than $c$ because the ledger
update propagates at one voxel per tick.  The entanglement entropy
of a subsystem satisfies an area law: $S \propto A/\ell_0^2$,
connecting to the Bekenstein--Hawking formula for black holes.
Decoherence is the loss of shared-entry correlations through
interaction with the environment.
\end{abstract}

%% ============================================================
\section{Introduction}
\label{sec:intro}
%% ============================================================

Since Einstein, Podolsky, and Rosen~\cite{EPR1935} and
Bell~\cite{Bell1964}, entanglement has been central to quantum
foundations.  Every interpretation of quantum
mechanics must account for entanglement: the existence of
correlations between spatially separated subsystems that cannot be
explained by any local hidden-variable theory.

RS~\cite{RS_Axioms} provides a clear ontological account: entanglement
is shared ledger structure.  No mystery, no spookiness---just two
subsystems sharing one constraint.

%% ============================================================
\section{Recognition Science Framework}
\label{sec:framework}
%% ============================================================

Recognition Science derives all physics from a single primitive,
the \emph{recognition cost functional} $J(x)$, uniquely determined
by three axioms.

\begin{definition}[RS Cost Functional]
For $x > 0$: $J(x) = \tfrac{1}{2}(x + x^{-1}) - 1$.
\end{definition}

\noindent\textbf{Axiom A1 (Normalization):} $J(1) = 0$.\\
\textbf{Axiom A2 (Recognition Composition Law, RCL):}
$J(xy) + J(x/y) = 2J(x)J(y) + 2J(x) + 2J(y)$ for all $x, y > 0$.\\
\textbf{Axiom A3 (Calibration):} $J''_{\log}(0) = 1$, where
$J_{\log}(t) := J(e^t)$.

\begin{theorem}[Cost Uniqueness]
\label{thm:unique}
The unique continuous solution to \textup{(A1)--(A3)} is
$J(x) = \tfrac{1}{2}(x + x^{-1}) - 1$.
\end{theorem}

\begin{proof}[Proof sketch]
Setting $f(t) = J_{\log}(t) + 1$ and rewriting \textup{(A2)} gives the
d'Alembert equation $f(s+t) + f(s-t) = 2f(s)f(t)$.  By the Acz\'el
classification theorem~\cite{Aczel1966}, the unique continuous solution
with $f(0) = 1$ and $f''(0) = 1$ is $f(t) = \cosh t$.
\end{proof}

\noindent The key property for entanglement is the reciprocal symmetry
$J(x) = J(x^{-1})$, which forces the double-entry structure: every
recognition event $x$ carries an equal-cost inverse $x^{-1}$.  Shared
double-entry pairs across two subsystems are the RS ontological basis of
entanglement.

%% ============================================================
\section{Entanglement from Shared Entries}
\label{sec:shared}
%% ============================================================

\begin{definition}[Entangled Subsystems]
\label{def:entangled}
Two subsystems $A$ and $B$ are entangled if and only if they
share at least one ledger entry---a single constraint that
correlates their outcomes.
\end{definition}

\begin{proposition}[Inseparability from Shared Entries]
\label{thm:inseparable}
If $A$ and $B$ share a ledger entry, then their joint state
cannot be written as a product:
\begin{equation}
  |\Psi_{AB}\rangle \neq |\psi_A\rangle \otimes |\psi_B\rangle.
\end{equation}
\end{proposition}

\begin{proof}[Structural argument]
A shared ledger entry constrains the outcomes of $A$ and $B$
simultaneously: if $A$ commits to outcome $|a_k\rangle$, the
shared entry fixes $B$'s outcome to $|b_k\rangle$.  A product
state $|\psi_A\rangle \otimes |\psi_B\rangle$ has the property
that $A$'s and $B$'s outcomes are statistically independent;
a shared entry enforces a correlation between them.  This correlation
cannot be expressed as a product of independent $A$ and $B$ states,
so $|\Psi_{AB}\rangle$ is inseparable.

\emph{Note}: This argument identifies shared ledger entries with
entangled states; the RS claim is that inseparability is a
\emph{definition} at the ledger level (a shared entry is by
definition a constraint binding both subsystems) rather than a
derived property of a pre-existing Hilbert space.
\end{proof}

\begin{corollary}[Entanglement Is Binary]
Two subsystems are either entangled (share at least one entry) or
separable (share no entries).  There is no ``partial'' sharing of
a single entry.
\end{corollary}

\begin{remark}
The \emph{degree} of entanglement (e.g., concurrence, entanglement
entropy) reflects the number and weight of shared entries, not a
continuum of sharing of a single entry.
\end{remark}

%% ============================================================
\section{Quantum Correlations}
\label{sec:correlations}
%% ============================================================

\begin{definition}[Quantum Correlation Function]
For a Bell pair with measurement settings $a$ and $b$:
\begin{equation}
  E(a, b) = -\cos(a - b).
\end{equation}
\end{definition}

\begin{proposition}[Bell Violation Verification --- Using Standard QM Correlations]
\label{thm:bell}
The CHSH combination
$S = E(a_1, b_1) - E(a_1, b_2) + E(a_2, b_1) + E(a_2, b_2)$
satisfies:
\begin{itemize}
  \item Classical bound (local hidden variables):
  $|S| \leq 2$.
  \item Quantum bound (shared entries):
  $|S| \leq 2\sqrt{2} \approx 2.83$ (Tsirelson bound~\cite{Tsirelson1980}).
\end{itemize}
For optimal measurement settings ($a_1 = 0$, $a_2 = \pi/2$,
$b_1 = \pi/4$, $b_2 = 3\pi/4$), the ledger's shared-entry
structure achieves $|S| = 2\sqrt{2}$, violating the classical bound.
\end{proposition}

\begin{proof}
We evaluate $S$ for optimal CHSH settings $a_1 = 0$, $a_2 = \pi/2$,
$b_1 = \pi/4$, $b_2 = 3\pi/4$, with $E(a,b) = -\cos(a - b)$:
\begin{align*}
  E(0,\tfrac{\pi}{4}) &= -\cos(-\tfrac{\pi}{4}) = -\tfrac{1}{\sqrt{2}},\\
  E(0,\tfrac{3\pi}{4}) &= -\cos(-\tfrac{3\pi}{4}) = +\tfrac{1}{\sqrt{2}},\\
  E(\tfrac{\pi}{2},\tfrac{\pi}{4}) &= -\cos(\tfrac{\pi}{4}) = -\tfrac{1}{\sqrt{2}},\\
  E(\tfrac{\pi}{2},\tfrac{3\pi}{4}) &= -\cos(-\tfrac{\pi}{4}) = -\tfrac{1}{\sqrt{2}}.
\end{align*}
Therefore:
\begin{equation}
  S = E(a_1,b_1) - E(a_1,b_2) + E(a_2,b_1) + E(a_2,b_2)
  = -4 \cdot \tfrac{1}{\sqrt{2}}
  = -2\sqrt{2}.
\end{equation}
So $|S| = 2\sqrt{2} > 2$, violating the classical bound.
The correlation $E(a,b) = -\cos(a-b)$ is not decomposable as
$\int \lambda(\xi)\,f_a(\xi)\,g_b(\xi)\,d\xi$ for any local
hidden-variable distribution $\lambda$, as shown by Bell's
theorem; the shared ledger entry provides the non-local constraint
that produces this correlation.
\end{proof}

\begin{remark}[Status of the Bell violation derivation]
\label{rem:bell-status}
The computation above verifies that the quantum correlation
function $E(a,b) = -\cos(a-b)$ achieves $|S| = 2\sqrt{2}$ for
the specified settings.  The underlying claim---that the RS
ledger's shared-entry mechanism \emph{produces} this correlation
function---rests on the identification of the shared entry with
the singlet state $|\Psi^-\rangle = (|01\rangle -
|10\rangle)/\sqrt{2}$.  The derivation of the singlet correlation
$E(a,b) = -\cos(a-b)$ directly from the RS double-entry structure
(rather than importing it from standard QM) is an identified open
problem.  The quantum-mechanical result is experimentally
established~\cite{Aspect1982}, and RS provides the ontological
account of what the shared entry \emph{is}; a first-principles
RS derivation of the correlation function is a future goal.
\end{remark}

%% ============================================================
\section{No Superluminal Signaling}
\label{sec:no-signaling}
%% ============================================================

\begin{proposition}[No-Signaling]
\label{thm:no-signal}
Entanglement does not enable faster-than-$c$ communication.
\end{proposition}

\begin{proof}
The shared constraint is established when the subsystems interact
(at speed $\leq c$).  Subsequent measurements reveal pre-existing
correlations but do not transmit new information.  Alice's
measurement yields $\pm 1$ with equal probability regardless of her
choice of axis; Bob's marginal distribution is unchanged:
\begin{equation}
  P_B(b) = \sum_a P(a, b) = \frac{1}{2},
\end{equation}
independent of Alice's setting.  The ledger update propagates at
$c = \ell_0/\tau_0$.
\end{proof}

%% ============================================================
\section{Entanglement Entropy}
\label{sec:entropy}
%% ============================================================

\begin{definition}[Entanglement Entropy]
The entanglement entropy of subsystem $A$ in the state
$|\Psi_{AB}\rangle$ is the von Neumann entropy of $A$'s reduced
density matrix:
\begin{equation}
  S_A = -\mathrm{Tr}(\rho_A \ln \rho_A),
  \quad \rho_A = \mathrm{Tr}_B |\Psi_{AB}\rangle
  \langle\Psi_{AB}|.
\end{equation}
\end{definition}

\begin{proposition}[Area Law --- Structural Argument]
\label{thm:area}
The entanglement entropy of a subsystem satisfies an area law:
\begin{equation}
  \boxed{S_A \leq \frac{A_{\partial}}{\ell_0^2}\ln\varphi,}
\end{equation}
where $A_{\partial}$ is the boundary area between $A$ and $B$
and $\ell_0$ is the voxel size.
\end{proposition}

\begin{proof}[Structural argument]
Shared ledger entries between $A$ and $B$ exist only at voxels
adjacent to the boundary $\partial(A|B)$, since entries deep inside
$A$ or $B$ are unconstrained by the opposing subsystem.
The number of boundary voxels scales as $N_{\partial} = A_{\partial}/\ell_0^2$.
In RS, the maximum entropy per voxel (from the $J$-cost binary
decision at each boundary site) is $\ln\varphi$ per $\varphi$-bit
cell, giving $S_A \leq N_{\partial}\ln\varphi$.
The proportionality $S_A \propto A_{\partial}/\ell_0^2$ follows.

\emph{Status note}: The identification of ``at most one $\varphi$-bit
per boundary voxel'' is a structural claim that requires a derivation
from the RS boundary Hamiltonian.  The area law itself is a well-known
result in quantum information and holography; the RS content is the
proposed mechanism (boundary voxels carrying $\varphi$-bits) and the
specific coefficient $\ln\varphi$.
\end{proof}

\begin{remark}[Coefficient in the area law]
\label{rem:area-coeff}
The coefficient in $S_A \leq (A_{\partial}/\ell_0^2)\ln\varphi$ is
the RS $\varphi$-bit entropy per cell.  The Bekenstein--Hawking
formula $S_{\mathrm{BH}} = A/(4\ell_0^2)\ln\varphi$ is a special
case with the factor of $1/4$ arising from the double-entry geometry
(see companion paper on BH entropy).  The derivation of the precise
factor of $1/4$ from the ledger's double-entry bookkeeping structure
is an identified open problem (treated in the companion paper).
\end{remark}

\begin{corollary}[Connection to Black Hole Entropy]
The Bekenstein--Hawking entropy $S_{\mathrm{BH}} = A/(4\ell_P^2)$
is a special case of the area law: a black hole is maximally
entangled with its exterior (every boundary voxel carries maximum
entropy).
\end{corollary}

%% ============================================================
\section{Decoherence}
\label{sec:decoherence}
%% ============================================================

\begin{proposition}[Decoherence from Environmental Sharing]
When a quantum system $A$ interacts with its environment $E$,
shared entries proliferate between $A$ and $E$.  The effective
state of $A$ becomes mixed:
\begin{equation}
  \rho_A = \mathrm{Tr}_E |\Psi_{AE}\rangle\langle\Psi_{AE}|
  \approx \sum_k p_k |a_k\rangle\langle a_k|.
\end{equation}
The interference (off-diagonal) terms are suppressed as the shared
entries with $E$ proliferate and become inaccessible.
\end{proposition}

\begin{remark}
Decoherence does not destroy entanglement; it distributes it into
the environment.  The total entanglement is conserved
(a structural consequence of double-entry bookkeeping: every
shared entry created between $A$ and~$E$ accounts for the
entanglement lost between $A$ and~$B$).  The \emph{accessible}
entanglement (measured by the subsystem entropy) decreases as
environmental sharing increases.
\end{remark}

%% ============================================================
\section{Comparison}
\label{sec:comparison}
%% ============================================================

\begin{center}
\renewcommand{\arraystretch}{1.3}
\begin{tabular}{@{}lccc@{}}
\toprule
& \textbf{Copenhagen} & \textbf{MWI} & \textbf{RS} \\
\midrule
Ontology of $|\Psi\rangle$ & Epistemic & Ontic & Computational \\
Entanglement & Correlation & Universal $\psi$ &
Shared entries \\
Bell violation & ``Shut up, calculate'' & Guaranteed &
Guaranteed \\
No-signaling & Postulated & From linearity &
From causality \\
Decoherence & Practical & Branching & Environmental sharing \\
\bottomrule
\end{tabular}
\end{center}

%% ============================================================
\section{Predictions}
\label{sec:predictions}
%% ============================================================

\begin{enumerate}
  \item \textbf{$|S| \leq 2\sqrt{2}$}: No experiment should
  exceed the Tsirelson bound.

  \item \textbf{Area law}: Entanglement entropy scales with
  boundary area, not volume.

  \item \textbf{No signaling}: All future experiments should
  confirm no superluminal communication via entanglement.

  \item \textbf{Falsifier}: $|S| > 2\sqrt{2}$ or superluminal
  signaling would falsify RS.
\end{enumerate}

%% ============================================================
\section{Conclusion}
\label{sec:conclusion}
%% ============================================================

Entanglement is shared ledger structure.  Two subsystems sharing a
ledger entry cannot be described by a product state.  Bell
violations are automatic; no-signaling is guaranteed by causality.
Entanglement entropy satisfies an area law, connecting to black
hole thermodynamics.  Decoherence is the proliferation of shared
entries with the environment.

%% ============================================================
\begin{thebibliography}{99}

\bibitem{Aczel1966}
J.~Acz\'el, \emph{Lectures on Functional Equations and Their Applications},
Academic Press, New York (1966).

\bibitem{RS_Axioms}
J.~Washburn, ``The Algebra of Reality: A Recognition Science Derivation
of Physical Law,'' \emph{Axioms} \textbf{15}(2), 90 (2026).
\texttt{doi:10.3390/axioms15020090}

\bibitem{Bell1964}
J.~S.~Bell, ``On the Einstein Podolsky Rosen Paradox,''
Physics \textbf{1}, 195 (1964).

\bibitem{EPR1935}
A.~Einstein, B.~Podolsky, and N.~Rosen, ``Can Quantum-Mechanical
Description of Physical Reality Be Considered Complete?''
Phys.\ Rev.\ \textbf{47}, 777 (1935).

\bibitem{Tsirelson1980}
B.~S.~Tsirelson, ``Quantum Generalizations of Bell's Inequality,''
Lett.\ Math.\ Phys.\ \textbf{4}, 93 (1980).

\bibitem{Aspect1982}
A.~Aspect, P.~Grangier, and G.~Roger, ``Experimental Realization
of Einstein-Podolsky-Rosen-Bohm Gedankenexperiment: A New
Violation of Bell's Inequalities,'' Phys.\ Rev.\ Lett.\
\textbf{49}, 91 (1982).

\end{thebibliography}

\end{document}
